Intuitively, \ommt is a declarative language for theories and views over an arbitrary object language.
% Its treatment of object languages is abstract enough to subsume most logics and type theories which are practically relevant.
For the purposes of this thesis, we will work with the (only slightly simplified) \hypertarget{mmt}{grammar given in \autoref{fig:prelim:mmt:grammar}}.

\begin{figure}[ht]%\centering%\vspace*{-1em}
  \begin{framed}
    \begin{tabular}{rl@{\quad}l|l}
      \cn{Thy} 	$::=$ 	& $T[:T]=\{ (\cn{Inc})^\ast\ (\cn{Const})^\ast \}$	& Theory				& \multirow{2}{*}{Modules} \\
      \cn{View} 	$::=$	& $T : \viewto TT = \{ (\cn{Ass})^\ast \}$			& View 				& \\ \hline
      \cn{Const} 	$::=$ 	& $C [\mmttp t] [\mmtdf t]$ 								& Constant Declarations& \multirow{3}{*}{Declarations} \\
      \cn{Inc}	$::=$	& $\incl T$								& Includes			& \\
      \cn{Ass}	$::=$	& $C \mmtass t$									& Assignments		& \\ \hline
      $\Gamma$	$::=$	& $(x [: t][:= t])^\ast$						& Variable Contexts	& \multirow{2}{*}{Objects} \\
      $t$ 		$::=$ 	& $x \mid C \mid \ombind{t}{\Gamma}{t} $		& Terms 	& \\%\hline\hline
%      $T$		$::=$ 	& \cn{String}								& Module Names 		& \multirow{3}{*}{Strings} \\
%      $C$		$::=$ 	& \cn{String}								& Constant Names		& \\
%      $x$		$::=$ 	& \cn{String}								& Variable Names		& \\
    \end{tabular}
    \\\ \\
    $T$ represents a module name, $C$ a constant name and $x$ a variable name
  \end{framed}%\vspace*{-.5em}
  \caption{The MMT Grammar}\label{fig:prelim:mmt:grammar}
\end{figure}

\subsubsection{Modules and Declarations}

%Fig.~\ref{fig:prelim:notations:mmt} gives an overview of the fundamental \mmt concepts.
In the simplest case, theories $T$ are lists of \textbf{constant declarations} $c\mmttp E$, where $E$ is an expression that may use the previously declared constants.
Naturally, $E$ must be subject to some type system (which MMT is also parametric in), for the purposes of this thesis primarily the logical framework \LF (see \autoref{sec:intro:prelim:lf}) and extensions thereof.
%We say that $\Sigma'$ includes $\Sigma$ if it contains every constant declaration of $\Sigma$.\medskip

\mmt achieves language-independence through the use of \textbf{meta-theories}: every \mmt-theory may designate a previously defined theory as its meta-theory.
For example, when we represent the HOL Light library in MMT, we first write a theory $L$ for the logical primitives of HOL Light.
Then each theory in the HOL Light library is represented as a theory with $L$ as its meta-theory.
In fact, we usually go one step further: $L$ itself is a theory, whose meta-theory is a logical framework such as \LF.
That allows $L$ to concisely define the syntax and inference system of HOL Light. \medskip

%\begin{figure}[htb]\centering
%\begin{tabular}{|l|l|l|}\hline
%\multicolumn{3}{|c|}{meta-theory: a fixed theory $M$}\\\hline
%       & Theory $\Sigma$    & View $\sigma:\Sigma\to\Sigma'$ \\\hline
%set of & typed constant declarations $c:E$ & assignments $c\mapsto E'$ \\
%$\Sigma$-expressions $E$ & formed from $M$- and $\Sigma$-constants  & mapped to $\Sigma'$ expressions \\\hline
%\end{tabular}
%\caption{Overview of MMT Concepts}\label{fig:prelim:notations:mmt}
%\end{figure}

Correspondingly, a \textbf{view} $V:\viewto T{T'}$ is a list of \textbf{assignments} $c \mmtass{e'}$ of $T'$-expressions $e'$ to $T$-constants $c$.
To be well-typed, $V$ must preserve typing, i.e., we must have $\vdash_{T'}e':\ov{V}(E)$.
Here $\ov{V}$ is the homomorphic extension of $V$, i.e., the map of $T$-expressions to $T'$-expressions that substitutes every occurrence of a $T$-constant with the $T'$-expression assigned by $V$.

We call $V$ \textbf{simple} if the expressions $e'$ are always $T'$-\emph{constants} rather than complex expressions.
The type-preservation condition for an assignment $c\mmtass c'$ reduces to $\ov{V}(E)\mmtass E'$ where $E$ and $E'$ are the types of $c$ and $c'$.
We call $V$ \textbf{partial} if it does not contain an assignment for every $T$-constant and \textbf{total} otherwise.
% A partial view from $T$ to $T'$ is the same as a total view from some theory included by $\Sigma$ to $\Sigma'$.

Importantly, we can then show generally at the \mmt-level that if $V$ is well-typed, then $\ov{V}$ preserves all typing and equality judgments over $T$.
In particular, if we represent proofs as typed terms (see \autoref{sec:intro:prelim:lf}), views preserve the theoremhood of propositions.
This property makes views so valuable for structuring, refactoring, and integrating large corpora.\medskip

Syntactically:
\begin{itemize}
	\item Theories $T: T' = \{ \incl{T_1}\ldots\incl{T_n}\quad C_1\ldots C_m \}$ have a module name $T$, an optional \emph{meta-theory} $T'$ and a body consisting of \emph{includes} of other theories $T_1\ldots T_n$
		 and a list of \emph{constant declarations} $C_1\ldots C_n$. A well-formed expression (see below) over $T$ is allowed to only reference constants that are declared in $T$, $T'$ or in one of the $T_i$ included.
	\item Constant declarations $C[\mmttp t][\mmtdf d]$ have a name $C$ and two optional term components; a \emph{type} ($[\mmttp t]$), and a \emph{definition} ($[\mmtdf d]$).
	\item Views $V : \viewto{T_1}{T_2} = \{ \ldots \}$ have a module name $V$, a domain theory $T_1$, a codomain theory $T_2$ and a body consisting of assignments $C \mmtass t$.
\end{itemize}

\begin{remark}\label{remark:prelim:mmt:syntax}
The module system is conservative: every theory can be elaborated (\emph{flattened}) into one that only declares constants, by recursively replacing every $\incl\cdot$ statement by the list of constants declared in the included theory. We will occasionally reference the result of flattening a theory $T$ as $T^\flat$.

Similarly, we can eliminate \emph{defined} constants $C\mmtdf d$ by replacing every occurence of the constant symbol $C$ by its definition $d$ (\emph{definition expansion}).

We allow definitions in variable contexts for purely technical reasons, most notably in the context of record types (see \autoref{sec:lf:records:def}). As a result, there is a notable similarity between variable contexts $\Gamma$ and theories; in fact, we can think of flattened theories as named variable contexts.
\end{remark}

% However, for our purposes, it suffices to say that the meta-theory is some fixed theory relative to which all concepts are defined.
% Thus, we assume that $\Sigma$ and $\Sigma'$ have the same meta-theory $M$, and that $\ov{\sigma}$ maps all $M$-constants to themselves.
%\begin{figure}% r{4.2cm}\vspace*{-2em}
%\begin{center}\begin{tabular}{l@{\quad}c@{\quad}l}
% $\Gamma$	& $::=$	& $(x:E)^\ast$ \\
% $E$      & $::=$ & $c \mid x \mid \ombind{E}{\Gamma}{E^+}$ \\
%\end{tabular}\end{center}%\vspace*{-2em}
%
%\caption{\mmt's Object Syntax}\label{fig:prelim:notations:expressions}
%\end{figure}

\subsubsection{Expressions}

It remains to define the exact syntax of expressions.   
In the grammar in \autoref{fig:prelim:mmt:grammar}, $C$ refers to constants  and $x$ refers to bound variables.

Simple expressions are either references $C$ to constants (of the meta-theory, an included theory or previously declared in the current theory) or to bound variables $x$.
Complex expressions are of the form $\ombind{o}{x_1:t_1,\ldots,x_m:t_m}{a_1,\ldots,a_n}$, where
\begin{compactitem}
 \item $o$ is the operator that forms the complex expression,
 \item $x_i:t_i$ declares a variable of type $t_i$ that is bound by $o$ in subsequent variable declarations and in the arguments,
 \item $a_i$ is an argument of $o$
\end{compactitem}
The bound variable context may be empty, and we write $\oma{o}{\vec{a}}$ instead of $\ombind{o}{\cdot}{\vec{a}}$.
For example, the axiom $\forall x:\cn{set},y:\cn{set}.\; \cn{finite}(x) \wedge y
\subseteq x \Rightarrow \cn{finite}(y)$ would instead be written as 
\[\expr{\ombind{\forall}{x:set,y:set}{\oma{\Rightarrow}{\oma{\wedge}{\oma{finite}{x},\oma{\subseteq}{y,x}},\oma{finite}{y}}}}\]

\subsubsection{URIs}

To refer to constants (and modules), \ommt employs globally unique URIs, which are composed of a namespace, the name of the (containing) module and the name of a constant, separated by question marks.

Hence, \mmt URIs are triples of the form 
 \[\mathtt{NAMESPACE} \;?\; \mathtt{MODULE} \;?\; \mathtt{SYMBOL}\]
 The namespace part is a URI that serves as a globally unique root identifier of a corpus,
 e.g. \url{http://mathhub.info/MyLogic/MyLibrary}. It is not necessary (although often
 useful) for namespaces to also be URLs, i.e., a reference to a physical location.  But
 even if they are URLs, we do not specify what resource dereferencing should return.  Note
 that because \mmt URIs use $?$ as a separator, $\mathtt{MODULE} \;?\; \mathtt{SYMBOL}$ is
 the query part of the URI, which makes it easy to implement dereferencing in practice.\footnote{For simplicity in the remaining part of the paper we will rarely use complete HTTP links, but rather use single keyword abbreviations.}

The module and symbol parts of an \mmt URI are logically meaningful names defined in the corpus: The module is the container (i.e. a theory, view or advanced structural feature) and the symbol is a name inside the module (of a constant, include or more advanced feature).
Both module and symbol name may consist of multiple /-separated segments to allow for nested modules and qualified symbol names. 

\mmt URIs allow arbitrary Unicode characters.  However, $?$ and $/$, which are used as
delimiters, as well as any character not legal in URIs must be escaped using the
\%-encoding (RFC 3986/7 standard).

By using URIs, namespaces have the great advantage of being guaranteed to be globally unique.
This comes at the prize of being rather long.
However, the CURIE standard can be used to introduce short prefixes.

\subsubsection{Structures}
\mmt offers an additional theory morphism that we will occasionally but rarely use in this thesis. A \emph{structure} $S:\viewto{T_1}{T_2}$ is a theory morphism that behaves like an include in that it makes all constants in $T_1$ accessible to $T_2$. Unlike includes, however, structures 
\begin{enumerate}
	\item can modify constants in the domain theory, e.g. by augmenting constants with new definientia, aliases and notations -- provided, the typing preservation condition for theory morphisms is maintained, and
	\item the symbols included via a structure are systematically renamed to be unique -- a constant $c$ in $T_1$ will be visible in $T_2$ as $S/c$. Consequently, unlike includes, structures are not idempotent.
\end{enumerate}
\begin{example}
	Assume we have a theory \cn{Monoid} (declaring a binary operation $\circ$ and a unit $e$), which is extended by a theory \cn{Group}. We can specify a theory \cn{Ring} by including two structures $\cn{plus}:\viewto{\cn{Group}}{\cn{Ring}}$ and $\cn{times}:\viewto{\cn{Monoid}}{\cn{Ring}}$. Unlike with theory includes, the two resulting morphisms from \cn{Monoid} to \cn{Ring} (one directly, and one via \cn{Group}) are \emph{not} identified, giving us two distinct operations. Furthermore, \cn{plus} and \cn{times} can rename the operations to $+$ and $\cdot$ respectively, and rename the units to 0 and 1.
\end{example}

\subsection{MathHub}\label{sec:intro:mathhub}
MMT is used as a backend for the \mathhub system~\cite{mathhub,MathHub:on} -- an online portal for mathematical documents. It is available on \url{mathhub.info} and hosts various \mmt repositories, interfaces for browsing and semantic search~\cite{mathwebsearch}, as well as several additional applications, such as the \emph{semantic glossary for mathematics (SMGLoM)}~\cite{GinIanJuc:spsttom16}.

Most of the specific \ommt formalizations in this thesis are hosted on MathHub's GitLab instance (\url{gl.mathhub.info}) and linked accordingly.


%\lstDeleteShortInline*

%An assignment in a view $V:T_1\to T_2$ is syntactically well-formed if for any assignment $C=t$ contained, $C$ is a constant declared in the flattened domain $T_1$ and $t$ is a syntactically well-formed term in the codomain $T_2$. We call a view \emph{total} if all \emph{undefined} constants in the domain have a corresponding assignment and \emph{partial} otherwise.

% Finally, we remark on a few additional features of the MMT language that are important for large-scale case studies but not critical to understand the basic intuitions of results.
% MMT provides a module system that allows theories to instantiate and import each other. The module system is conservative: every theory can be \emph{elaborated} into one that only declares constants.
% MMT constants may carry an optional definiens, in which case we write $c:E=e$.
% Defined constants can be eliminated by definition expansion.

%\begin{oldpart}{FR: replaced with the above}
%
%In more detail:
%\begin{itemize}
%	\item Theories have a module name, an optional \emph{meta-theory} and a body consisting of \emph{includes} of other theories
%		 and a list of \emph{constant declarations}.
%	\item Constant declarations have a name and two optional term components; a \emph{type} ($[:t]$), and a \emph{definition} ($[=t]$).
%	\item Views $V : T_1 \to T_2$ have a module name, a domain theory $T_1$, a codomain theory $T_2$ and a body consisting of assignments
%		$C = t$.
%	\item Terms are either 
%		\begin{itemize}
%			\item variables $x$, 
%			\item symbol references $T?C$ (referencing the constant $C$ in theory $T$), 
%			\item applications $\oma{f}{a_1,\ldots,a_n}$ of a term $f$ to a list of arguments $a_1,\ldots,a_n$ or
%			\item binding application $\ombind{f}{x_1[:t_1][=d_1],\ldots,x_n[:t_n][=d_n]}{b}$, where $f$ \emph{binds} the variables
%				$x_1,\ldots,x_n$ in the body $b$ (representing binders such as quantifiers, lambdas, dependent type constructors etc.).
%		\end{itemize}
%\end{itemize}
%The term components of a constant in a theory $T$ may only contain symbol references to constants declared previously in $T$, or that are declared in some theory $T'$ (recursively) included in $T$ (or its meta-theory, which we consider an \emph{include} as well). 
%We can eliminate all includes in a theory $T$ by simply copying over the constant declarations in the included theories; we call this process \emph{flattening}. We will often and without loss of generality assume a theory to be \emph{flat} for convenience.
%
%An assignment in a view $V:T_1\to T_2$ is syntactically well-formed if for any assignment $C=t$ contained, $C$ is a constant declared in the flattened domain $T_1$ and $t$ is a syntactically well-formed term in the codomain $T_2$. We call a view \emph{total} if all \emph{undefined} constants in the domain have a corresponding assignment and \emph{partial} otherwise.
%
%\end{oldpart}


%\begin{example}[$\lambda$-calculus]\label{ex:lamsyn}
%To give an urtheory for the simply-typed $\lambda$-calculus in \mmt, we use $4$ constructors $\mathcal{C}=\{\type,\to,\lambda,@\}$.
%Alternatively, we can add constructors $\kind$ and $\Pi$ to obtain an urtheory for the dependently-typed $\lambda$-calculus.
%
%We introduce the usual notations for them as follows:
%
%\begin{center}
%\begin{tabular}{|l|l|l|}
%\hline
%Constructor       & Abstract \mmt expression     & Concrete notation\\
%\hline
%universe of types & $\cons{\type}{\cdot}{\cdot}$ & $\type$ \\
%function types    & $\cons{\to}{\cdot}{A,B}$     & $A\to B$  \\
%abstraction       & $\cons{\lambda}{x:A}{t}$     & $\lambda x:A.t$ \\
%application       & $\cons{@}{\cdot}{f,t}$       & $f\,t$\\
%\hline
%universe of kinds & $\cons{\kind}{\cdot}{\cdot}$ & $\kind$ \\
%dependent function types & $\Pi(x:A;t)$  & $\Pi x:A.t$ \\
%\hline
%\end{tabular}
%\end{center}
%\end{example}


%%% Local Variables:
%%% mode: latex
%%% eval:  (set-fill-column 5000)
%%% eval:  (visual-line-mode)
%%% TeX-master: "paper"
%%% End:

%  LocalWords:  fig:mmt textbf textbf mapsto vdash_ emph hline ctabular ldots,x_m:t_m vec
%  LocalWords:  ldots,a_n compactitem cdot vec ednote ex:lamsyn urtheory mathcal f,t y,x
%  LocalWords:  x:A;t oldpart fig:mmtgrammar centering vspace mdframed multirow ldots,x_n
%  LocalWords:  OAFproject:on vdash formalized formalizations tikzpicture ldots Arith s,t
%  LocalWords:  newpart forall subseteq Rightarrow forall Rightarrow subseteq t,s lf
